% Reviewer 4; OpenReview ID: hTMj
% Original review text. Author responses have not been drafted.
\section{Reviewer 4 (\texttt{hTMj})}
\label{reviewer:4}
\noindent\textbf{Rating:} 2: Strong rejection\par
\noindent\textbf{Confidence:} 4: The reviewer is confident but not absolutely certain that the evaluation is correct

\subsection*{Paper Summary}
This paper proposes R2-OPD, a selective on-policy distillation method motivated by the potential mismatch between teacher agreement and reasoning progress. The method samples continuations from reasoning prefixes to estimate changes in the probability of obtaining a correct answer, merges adjacent spans with sign-consistent progress, and masks supervision where progress and average distillation-loss rankings conflict. Overall, the paper proposes a highly incremental design, and the scope of the experiments, the analysis, and the results are all quite limited.

\subsection*{Strengths}
\begin{enumerate}
\item Teacher agreement can differ from reasoning progress, making selective distillation a useful direction. Applying existing progress-based supervision [2] to span-level filtering is a concrete proposal.
\item Table 1 reports improvements over OPD in avg@4 and pass@4. The masking and merging ablations demonstrate sensitivity to these design choices, without establishing the claimed mechanism.
\end{enumerate}

\subsection*{Weaknesses}
\begin{enumerate}
\item The method combines existing ideas from progress-based supervision and reward-guided distillation [1,2]. Its main distinction is using estimated progress to select individual reasoning spans, which appears to be a relatively narrow methodological extension.
\item Both evaluated students are small Qwen-derived models. This makes it difficult to tell whether the gains extend beyond these particular models. Results on another model family, such as Llama, and larger students would be needed to support broader claims. All evaluations use a single response-length limit. The paper does not show whether the gains remain when longer reasoning is allowed. RG-OPD is closely related to the proposed method but is not included in the experiments [1]. Comparisons with the closest existing methods and simpler reward-based filters are needed.
\item The filter relies on progress estimates from continuations limited to 300 tokens. It is unclear whether these short, answer-eliciting completions reliably reflect how much a reasoning step contributes to solving the problem. The paper does not examine how these estimates change with the continuation length or number of samples. Estimating progress requires additional generations from multiple prefixes of each eligible trajectory. The paper does not report the full cost of this procedure, so training each method for the same number of epochs does not establish equal resource use.
\item The theory explains how aggregation can cancel some estimation errors and reduce variance under its assumptions. It does not address how using noisy progress scores to choose the segments affects these results.
\item The results are reported without run-to-run variation or uncertainty estimates. This is a concern because improvements are modest for the second model pair and are not uniform across benchmarks. Repeated runs are needed to distinguish a reliable gain from variation in training and evaluation.
\end{enumerate}

[1] Akhondzadeh et al. Reward-Gated On-Policy Distillation. arXiv:2607.04037v1, 2026; Section 3.

[2] Setlur et al. Rewarding Progress: Scaling Automated Process Verifiers for LLM Reasoning. ICLR 2025; arXiv:2410.08146, Section 3.1.

[3] Agentica. DeepScaleR-Preview-Dataset.

\subsection*{Justification of Recommendation}
See the weaknesses.

\subsection*{Specific Points of Feedback for Rebuttal}
\begin{enumerate}
\item What benefit is specific to the proposed selection rule, compared with simpler ways of using process or outcome rewards?
\item Do the improvements hold for larger students and non-Qwen backbones, such as Llama?Does the advantage remain when all methods are allowed longer responses?
\item How does R2-OPD compare with RG-OPD and the other relevant baselines on both model pairs?
\item How reliable are the progress estimates when the continuation length or number of samples is changed?
\item How does selecting segments from the same noisy scores affect the theoretical argument?
\end{enumerate}

% ------------------------------------------------------------------
% Author response: leave blank until the collaborative drafting stage.
% Uncomment the following heading when ready to write a response.
% \subsection*{Author Response}
% TODO: Draft the response to this reviewer here.
